\section{Equivariant cohomology with supports}
\label{section-equivariant-cohomology-supports}

\noindent
Throughout this section a $G$-sheaf means an abelian $G$-sheaf. No extension
to modules over a general equivariant sheaf of rings is asserted.

\begin{definition}
\label{definition-equivariant-cohomology-family-supports}
\begin{reference}
\cite[Chapter V, \S 5.7, p. 219]{Tohoku}
\end{reference}
Let $G$ act on a topological space $X$, and let $\Phi$ be a family of
supports on $X$ in the sense of Definition
\ref{definition-family-supports}. For an abelian $G$-sheaf $\mathcal A$ on
$X$, set
$$
\Gamma^G_\Phi(X,\mathcal A)
=
\{s\in\Gamma(X,\mathcal A)^G\mid
\text{Supp}_X(s)\in\Phi\}.
$$
Thus the right hand side means the intersection of the invariant global
sections with the sections supported in $\Phi$; it does not require
$\Gamma_\Phi(X,\mathcal A)$ to be stable under $G$.

This defines an additive left exact functor
$$
\Gamma^G_\Phi(X,-):
\textit{Mod}_G(\underline{\mathbf Z}_X)
\longrightarrow
\textit{Ab};
$$
see Lemma \ref{lemma-equivariant-family-supports-invariant-saturation}.
The source category has enough injectives by Injectives, Lemma
\ref{injectives-lemma-equivariant-modules-enough-injectives}. We define
$$
H^n_\Phi(X;G,\mathcal A)
=
R^n\Gamma^G_\Phi(X,-)(\mathcal A),
\qquad n\geq 0,
$$
and put these groups equal to zero in negative degrees. If $\Phi$ is the
family of all closed subsets of $X$, this recovers
$H^n(X;G,\mathcal A)$ of Definition
\ref{definition-orbit-space-equivariant-cohomology}.
\end{definition}

\begin{lemma}[Invariant saturation of a family of supports]
\label{lemma-equivariant-family-supports-invariant-saturation}
\begin{reference}
\cite[Chapter V, \S 5.7, p. 219]{Tohoku}
\end{reference}
Let $G$ act on $X$, and let $\Phi$ be a family of supports on $X$. Put
$$
\Phi^G=
\left\{
Z\subset X\ \middle|\
\begin{aligned}
&Z\text{ is closed, and there is }W\in\Phi\text{ with }Z\subset W,\\
&gW=W\text{ for every }g\in G
\end{aligned}
\right\}.
$$
Then $\Phi^G$ is a family of supports, it is stable under the action of
$G$, and $\Phi^G\subset\Phi$. For every abelian $G$-sheaf $\mathcal A$
there is a functorial equality
$$
\Gamma^G_\Phi(X,\mathcal A)
=
\Gamma^G_{\Phi^G}(X,\mathcal A).
$$
For the family on the right, supported global sections form a
$\mathbf Z[G]$-module and
$$
\Gamma^G_{\Phi^G}(X,-)
=
(-)^G\circ\Gamma_{\Phi^G}(X,-).
$$
In particular, $\Gamma^G_\Phi(X,-)$ is additive and left exact.
\end{lemma}

\begin{proof}
Every family of supports contains the empty set. Hence the empty set is in
$\Phi^G$. If $Z_i\subset W_i$ exhibits $Z_i\in\Phi^G$ for $i=1,2$, then
$W_1\cup W_2$ is an invariant member of $\Phi$ containing
$Z_1\cup Z_2$. Downward closure is immediate. This proves that $\Phi^G$
is a family of supports. If $Z\subset W$ is such an exhibition, then
$gZ\subset W$ for every $g\in G$, so the family is $G$-stable. Downward
closure of $\Phi$ also gives $\Phi^G\subset\Phi$.

Let $s\in\Gamma(X,\mathcal A)^G$. The action isomorphisms on the stalks
carry $s_x$ to $s_{gx}$. Thus
$$
g\text{Supp}_X(s)=\text{Supp}_X(s)
$$
for every $g\in G$. If $\text{Supp}_X(s)\in\Phi$, this support itself is
an invariant member of $\Phi$, and hence it lies in $\Phi^G$. The converse
follows from $\Phi^G\subset\Phi$. This proves the functorial equality.

Because $\Phi^G$ is $G$-stable, the group action on global sections
preserves $\Gamma_{\Phi^G}(X,\mathcal A)$. This gives the last displayed
factorization. The functor $\Gamma_{\Phi^G}(X,-)$ is additive and left
exact by Lemma \ref{lemma-cohomology-family-supports}, and taking
$G$-invariants is additive and left exact. The final assertion follows.
\end{proof}

\begin{lemma}[Equivariant cohomology with supports on the orbit space]
\label{lemma-orbit-space-equivariant-support-spectral-sequence}
\begin{reference}
\cite[Chapter V, \S 5.7, pp. 219--220]{Tohoku}
\end{reference}
\begin{history}
The support subscript on the abutment in the cited source is essential. It
is the subscript restored in corrected-source decision
\texttt{TOH-SRC-P220-004}; omitting it would contradict the degree-zero
functor identity proved below.
\end{history}
Let $G$ act on $X$, let $Y=X/G$, and let $q:X\to Y$ be the orbit map. Let
$\Psi$ be a family of supports on $Y$, and put
$$
q^\dagger\Psi=
\{Z\subset X\mid Z\text{ is closed and }
\overline{q(Z)}\in\Psi\}.
$$
The closure in this formula is part of the definition; the image under an
arbitrary orbit map of a closed subset need not be closed. Then:
\begin{enumerate}
\item $q^\dagger\Psi$ is a family of supports on $X$, and it equals its
invariant saturation $(q^\dagger\Psi)^G$.
\item If $\mathcal A$ is an abelian $G$-sheaf on $X$ and an invariant
section $s\in\Gamma(X,\mathcal A)^G$ corresponds to
$\widetilde s\in\Gamma(Y,q_*^G\mathcal A)$, then
$$
\text{Supp}_Y(\widetilde s)
=
\overline{q(\text{Supp}_X(s))}.
$$
Consequently there is a natural equality of left exact functors
$$
\Gamma^G_{q^\dagger\Psi}(X,-)
=
\Gamma_\Psi(Y,-)\circ q_*^G.
$$
\item For every abelian $G$-sheaf $\mathcal A$ there is a functorial
first-quadrant spectral sequence
$$
E_2^{p,q}
=
H^p_\Psi\bigl(Y,\mathcal H^q(G,\mathcal A)\bigr)
\Longrightarrow
H^{p+q}_{q^\dagger\Psi}(X;G,\mathcal A).
$$
It is additive and functorial from the $E_2$-page onward and gives a
finite filtration on every abutment group.
\end{enumerate}
\end{lemma}

\begin{proof}
The empty set belongs to $q^\dagger\Psi$. If $Z'\subset Z$ are closed and
$Z\in q^\dagger\Psi$, then
$$
\overline{q(Z')}\subset\overline{q(Z)},
$$
so downward closure of $\Psi$ gives $Z'\in q^\dagger\Psi$. For two closed
subsets $Z_1,Z_2\subset X$, we have
$$
\overline{q(Z_1\cup Z_2)}
=
\overline{q(Z_1)}\cup\overline{q(Z_2)}.
$$
Thus closure under finite unions follows from the corresponding property
of $\Psi$. This proves that $q^\dagger\Psi$ is a family of supports.

Let $Z\in q^\dagger\Psi$, put $C=\overline{q(Z)}$, and set
$$
W=q^{-1}(C).
$$
The subset $W$ is closed, $G$-invariant, and contains $Z$. Since $q$ is
surjective, $q(W)=C$, and hence $W\in q^\dagger\Psi$. Therefore every
member of $q^\dagger\Psi$ is contained in an invariant member of the same
family. Lemma
\ref{lemma-equivariant-family-supports-invariant-saturation} gives
$(q^\dagger\Psi)^G=q^\dagger\Psi$.

We prove the support identity. Put $S=\text{Supp}_X(s)$. For $y\in Y$, the
definition of a stalk and of invariant direct image gives
\begin{align*}
\widetilde s_y=0
&\Longleftrightarrow
\text{there is an open }V\ni y\text{ such that }
s|_{q^{-1}(V)}=0\\
&\Longleftrightarrow
\text{there is an open }V\ni y\text{ such that }
V\cap q(S)=\emptyset\\
&\Longleftrightarrow
y\notin\overline{q(S)}.
\end{align*}
Hence $\text{Supp}_Y(\widetilde s)=\overline{q(S)}$. By Definition
\ref{definition-family-supports}, the section $\widetilde s$ is supported
in $\Psi$ if and only if $s$ is supported in $q^\dagger\Psi$. This proves
the functor identity.

Set
$$
F=q_*^G:
\textit{Mod}_G(\underline{\mathbf Z}_X)
\longrightarrow\textit{Ab}(Y)
\quad\text{and}\quad
T=\Gamma_\Psi(Y,-):\textit{Ab}(Y)\longrightarrow\textit{Ab}.
$$
These are additive left exact functors by Modules, Lemma
\ref{modules-lemma-invariant-direct-image-G-O-modules} and Lemma
\ref{lemma-cohomology-family-supports}. If $\mathcal I$ is an injective
abelian $G$-sheaf on $X$, then $F(\mathcal I)=q_*^G\mathcal I$ is flasque
by Injectives, Lemma
\ref{injectives-lemma-equivariant-injective-sections-invariants}. It is
therefore right acyclic for $T$ by Lemma
\ref{lemma-family-supports-flasque-soft-acyclic}.

The Grothendieck spectral sequence, Derived Categories, Lemma
\ref{derived-lemma-grothendieck-spectral-sequence}, now applies to
$T\circ F$. Definition
\ref{definition-orbit-space-equivariant-cohomology} identifies
$R^qF(\mathcal A)$ with $\mathcal H^q(G,\mathcal A)$, while Lemma
\ref{lemma-cohomology-family-supports} identifies $R^pT$ with
$H^p_\Psi(Y,-)$. The degree-zero functor identity above and Definition
\ref{definition-equivariant-cohomology-family-supports} identify the
derived composite with
$H^{p+q}_{q^\dagger\Psi}(X;G,\mathcal A)$. This proves (3).
\end{proof}

\begin{remark}
\label{remark-equivariant-support-second-spectral-sequence}
\begin{reference}
\cite[Chapter V, \S 5.7, p. 220]{Tohoku}
\end{reference}
For the invariant family $q^\dagger\Psi$, there is also a degree-zero
factorization
$$
\Gamma^G_{q^\dagger\Psi}(X,-)
=
(-)^G\circ\Gamma_{q^\dagger\Psi}(X,-).
$$
This factorization alone does not produce a second Grothendieck spectral
sequence. One would have to prove, in particular, that
$\Gamma_{q^\dagger\Psi}(X,\mathcal I)$ is right acyclic for $(-)^G$ for
every injective abelian $G$-sheaf $\mathcal I$; the preceding results do not
give that assertion. This is only a failure of the present hypotheses to
verify the Grothendieck acyclicity condition; it is not a nonexistence
claim. Thus we do not assert a spectral sequence with a putative initial page
$$
H^p\bigl(G,H^q_{q^\dagger\Psi}(X,\mathcal A)\bigr).
$$
The cited source next assumes, exactly, that every member of $\Psi$ has a
neighbourhood which is again a member of $\Psi$. Under that hypothesis it
passes through extension by zero and a filtered system of neighbourhoods in
order to construct a different second spectral sequence. We record the
hypothesis here because it marks the end of the obstruction statement, but
we do not yet identify supported cohomology with that filtered colimit,
construct the resulting spectral sequence, or interchange group cohomology
with a filtered colimit. Those are separate assertions requiring the
additional work carried out after the cited passage.
\end{remark}
